% =====================================================================
%  papers/shared/primer.tex
%
%  PLACEHOLDER created by T-005 so the build scaffold compiles end to end.
%  T-004 OWNS THIS FILE and will replace it wholesale.  Nothing here is
%  intended as final prose.
%
%  Contract this stub establishes, which T-004 should preserve:
%    * The file is \input inside a \section, so it starts at \subsection
%      level.  It must not issue \section or \chapter.
%    * It must not assume either paper's specific context (D-006).
%    * Every term of art gets \defterm{key}{printed term} at its defining
%      occurrence, and a matching \gentry line in shared/glossary.tex.
% =====================================================================

\subsection{What this primer assumes}
\label{sec:primer-assumes}

This primer assumes a reader who is comfortable with algebra and with reading a
graph, and who has no background in statistics whatsoever. Every term of art is
defined before it is used. Terms are set in \textbf{\textit{bold italic}} at the
point of definition and collected in the glossary at the end of the paper, each
with the section number where its definition lives.

\subsection{Locations, values, and neighbours}
\label{sec:primer-weights}

Suppose we have measured one number at each of $n$ places --- a rate, a count, a
score. Write $y_i$ for the value at place $i$. Spatial statistics asks whether
the arrangement of those numbers across the map carries information that the
numbers alone do not.

To ask that question we must first say which places count as neighbours of which
others. That bookkeeping is done by a \defterm{spatial-weights-matrix}{spatial
weights matrix}: a square table $\Wmat$ with one row and one column per place,
whose entry $w_{ij}$ records how strongly place $j$ counts as a neighbour of
place $i$. A common choice sets $w_{ij}=1$ when $j$ is one of the $k$ nearest
places to $i$ and $w_{ij}=0$ otherwise.

Reading the table as a map of who-is-connected-to-whom, a group of places that
can all reach one another by hopping along nonzero weights, and which no weight
connects to anything outside the group, is a \defterm{connected-component}{connected
component} of the neighbour graph. A neighbour graph with more than one component
has been cut into pieces that never exchange information.

\subsection{Evidence against a hypothesis}
\label{sec:primer-pvalue}

Statistical tests work by contradiction. We name a dull explanation for what we
observed, then ask how surprising the observation would be if that dull
explanation were true. The dull explanation is the
\defterm{null-hypothesis}{null hypothesis}; here it is usually ``the values were
scattered across the map with no regard to location.''

The measure of surprise is the \defterm{p-value}{p-value}: the probability of
seeing a result at least as extreme as the one observed, computed under the
assumption that the null hypothesis holds. A small p-value means the observation
sits far out in the tail of what the dull explanation predicts. It does not mean
the dull explanation is false, and it is not the probability that the dull
explanation is false.

\begin{keyinsight}[Placeholder]
This is placeholder primer content supplied by the build scaffold (T-005) so
that the document compiles, the glossary resolves, and the cross-reference
machinery can be checked. T-004 replaces this file with the real
from-first-principles exposition required by AC-003.
\end{keyinsight}
